A crescente demanda por comunicações sem fio impõe desafios significativos relacionados à interferência eletromagnética e à alocação eficiente do espectro de frequências. O problema da Rotulação–L(2,1) em grafos modela a atribuição de canais de forma a evitar interferências entre transmissores próximos, exigindo que vértices adjacentes tenham rótulos com diferença absoluta mínima de pelo menos duas unidades e vértices à distância 2 tenham diferença absoluta mínima de pelo menos uma unidade. Assim, uma rotulação-L(2,1) de um grafo $G$ é uma função $f$ que atribui inteiros não negativos ao conjunto $V(G)$ tal que $|f(u) - f(v)| \geq 2$ se $d(u, v) = 1$ e $|f(u) - f(v)| \geq 1$ se $d(u, v) = 2$. O \textit{span} de uma rotulação $f$ é o maior valor atribuído por $f$ a um vértice de $G$. O objetivo do problema de Rotulação–L(2,1) é encontrar uma rotulação válida $f$ que minimize esse valor. Contudo, encontrar essa rotulação é um problema NP-completo, o que inviabiliza abordagens exatas para instâncias de tamanho moderado ou grande. Neste contexto, este trabalho propõe a investigação do uso de Algoritmos Genéticos como alternativa heurística eficiente para a resolução do problema. Este trabalho tem como objetivo comparar a eficácia de diferentes variações de Algoritmos Genéticos aplicados ao Problema da Rotulação–L(2,1). A metodologia inclui a implementação dos algoritmos, aplicação em diferentes classes de grafos, calibração de parâmetros, e comparação com uma formulação de Programação Linear Inteira. Os resultados serão analisados com base na qualidade da solução, tempo de execução e estabilidade. Esta abordagem visa contribuir para a otimização da alocação espectral, com impactos relevantes na eficiência de redes de comunicação.

\palavraschave{rotulação–L(2,1); grafos; algoritmos genéticos; otimização; heurística.}